\documentclass[aps,preprint]{revtex4}%
\usepackage{amsfonts}
\usepackage{subfigure}
\usepackage{amsmath}
\usepackage{amssymb}
\usepackage{graphicx}%
\setcounter{MaxMatrixCols}{30}
%TCIDATA{OutputFilter=latex2.dll}
%TCIDATA{Version=5.50.0.2953}
%TCIDATA{CSTFile=revtex4.cst}
%TCIDATA{Created=Wednesday, May 30, 2007 12:29:48}
%TCIDATA{LastRevised=Wednesday, May 30, 2007 15:37:20}
%TCIDATA{<META NAME="GraphicsSave" CONTENT="32">}
%TCIDATA{<META NAME="SaveForMode" CONTENT="1">}
%TCIDATA{BibliographyScheme=Manual}
%TCIDATA{<META NAME="DocumentShell" CONTENT="Articles\SW\REVTeX 4">}
%BeginMSIPreambleData
\providecommand{\U}[1]{\protect\rule{.1in}{.1in}}
%EndMSIPreambleData
\newtheorem{theorem}{Theorem}
\newtheorem{acknowledgement}[theorem]{Acknowledgement}
\newtheorem{algorithm}[theorem]{Algorithm}
\newtheorem{axiom}[theorem]{Axiom}
\newtheorem{claim}[theorem]{Claim}
\newtheorem{conclusion}[theorem]{Conclusion}
\newtheorem{condition}[theorem]{Condition}
\newtheorem{conjecture}[theorem]{Conjecture}
\newtheorem{corollary}[theorem]{Corollary}
\newtheorem{criterion}[theorem]{Criterion}
\newtheorem{definition}[theorem]{Definition}
\newtheorem{example}[theorem]{Example}
\newtheorem{exercise}[theorem]{Exercise}
\newtheorem{lemma}[theorem]{Lemma}
\newtheorem{notation}[theorem]{Notation}
\newtheorem{problem}[theorem]{Problem}
\newtheorem{proposition}[theorem]{Proposition}
\newtheorem{remark}[theorem]{Remark}
\newtheorem{solution}[theorem]{Solution}
\newtheorem{summary}[theorem]{Summary}
\newenvironment{proof}[1][Proof]{\noindent\textbf{#1.} }{\ \rule{0.5em}{0.5em}}
\begin{document}
\preprint{HEP/123-qed}
\title[Short title for running header]{Long title that appears on the title page}
\author{First Author}
\affiliation{My Institution}
\author{Second Author}
\affiliation{My Institution}
\author{Third Author}
\affiliation{Other Institution}
\keywords{one two three}
\pacs{PACS number}

\begin{abstract}
Shell document for REV\TeX{} 4.

\end{abstract}
\volumeyear{year}
\volumenumber{number}
\issuenumber{number}
\eid{identifier}
\date[Date text]{date}
\received[Received text]{date}

\revised[Revised text]{date}

\accepted[Accepted text]{date}

\published[Published text]{date}

\startpage{101}
\endpage{102}
\maketitle
\tableofcontents

In fig's \ref{fig:figone-first}, \ref{fig:figone-second}, \ref{fig:figone-third}, figs \ref{fig:figtwo-first} and   \ref{fig:figtwo-second}%

\begin{figure}
\begin{center}
\subfigure{\label{fig:figone-first}\includegraphics[width=85mm]{transmitted.eps}}
\subfigure{\label{fig:figone-second}\includegraphics[width=85mm]{absorbed.eps}}
\subfigure{\label{fig:figone-third}\includegraphics[width=85mm]{reflected.eps}}

\end{center}
\\
\caption{(Color online). Plots of the magnitudes of (a) the power transmitted,
(b) absorbed and (c) reflected by the light incident on the slab formed by a
LC that has relative permittivity of 2,3 and 4.}

\label{fig:figone}
\end{figure}

\begin{figure}
\begin{center}
\subfigure{\label{fig:figtwo-first}\includegraphics[width=85mm]{indexofrefract.eps}}
\subfigure{\label{fig:figtwo-second}\includegraphics[width=85mm]{impedance.eps}}

\end{center}
\\
\caption{(Color online). Plots of (a) the real (solid)
and imaginary (dashed) parts of the bulk
refractive index and (b) of the real (solid)
and imaginary (dashed) parts of the intrinsic
impedance of the slabs formed by a LC that is tuned to have a relative
permittivity of 2,3 and 4.}

\label{fig:figtwo}
\end{figure}


\end{document} 